% ============================================================
% High-Density Object Segmentation with Soft-NMS
% IEEE Conference Paper — Phase 2 Final Report
% ============================================================
\documentclass[conference]{IEEEtran}

% ---- Packages ----
\usepackage[utf8]{inputenc}
\usepackage[T1]{fontenc}
\usepackage{amsmath,amssymb}
\usepackage{graphicx}
\usepackage{booktabs}
\usepackage{hyperref}
\usepackage{cite}
\usepackage{xcolor}
\usepackage{subcaption}
\usepackage{algorithm}
\usepackage{algorithmic}
\usepackage{multirow}

% ---- Custom commands ----
\newcommand{\eg}{\textit{e.g.}}
\newcommand{\ie}{\textit{i.e.}}

% ---- Document ----
\begin{document}

\title{High-Density Object Segmentation with Soft-NMS:\\
Instance-Level Detection of Heavily Overlapping Objects}

\author{
  \IEEEauthorblockN{Samarth Shekhar \and Sanskar}
  \IEEEauthorblockA{Bell Labs Research Project\\
  Rishihood University}
}

\maketitle

% ============================================================
\begin{abstract}
Many real-world vision systems operate in environments where objects are
densely packed and heavily occluded. Standard Non-Maximum Suppression (NMS)
aggressively discards valid detections in such settings, leading to
significant under-counting. This paper presents a systematic study of
instance-level detection and segmentation for 1--50 heavily overlapping
objects in static images. We implement a progressive four-level pipeline
spanning heuristic baselines, classical computer vision methods, a
Random Forest machine learning classifier, and deep learning detectors
augmented with Soft-NMS. We further introduce a lightweight density
estimation head and evaluate performance on a filtered subset of SKU-110K
and a synthetic overlapping-shapes dataset. Our results demonstrate that
replacing Hard NMS with Gaussian Soft-NMS yields consistent improvements:
\textbf{3.4\%} mAP@0.5 gain and \textbf{57\%} reduction in counting
error (MAE 4.2 $\to$ 1.8) over Hard NMS under extreme overlap scenarios.
Robustness analysis under Gaussian noise, blur, and brightness shifts
confirms that performance degrades gracefully across all perturbation
conditions, with mAP dropping by no more than 4.1\% under the harshest
corruption tested.
\end{abstract}

\begin{IEEEkeywords}
Object detection, instance segmentation, Soft-NMS, dense scenes, occlusion,
SKU-110K, retail shelf monitoring, Random Forest, robustness
\end{IEEEkeywords}

% ============================================================
\section{Introduction \& Problem Statement}
\label{sec:introduction}

Object detection has achieved remarkable progress on standard benchmarks,
yet most systems implicitly assume that objects are well-separated in the
image. In many real-world scenarios---retail shelf monitoring, warehouse
inventory counting, crowd analysis, and biological cell imaging---objects
are \emph{densely packed and heavily occluded}, violating this assumption.

This project targets \textbf{instance-level detection and segmentation of
1--50 heavily overlapping objects} in static images. In dense retail shelves,
products may overlap by 50\% or more of their bounding box area. In crowd
scenes, pedestrians frequently share significant spatial extent. Standard
Non-Maximum Suppression (NMS) treats overlapping detections as duplicates
and suppresses all but the highest-scoring box, leading to catastrophic
under-detection when genuinely distinct objects share spatial extent.

\textbf{Key contributions.} This work makes the following contributions:
\begin{enumerate}
  \item We implement and compare four progressively complex detection
    levels: heuristic baseline, advanced classical CV, Random Forest ML
    classification, and deep learning with Soft-NMS.
  \item We introduce a \emph{configurable Soft-NMS module} supporting
    Gaussian, Linear, and Hard decay modes as a drop-in replacement for
    standard NMS in Faster~R-CNN.
  \item We add a \emph{coarse density estimation head} that predicts
    approximate object counts from FPN features, acting as both a
    training regularizer and a fast inference proxy.
  \item We conduct comprehensive robustness analysis under five
    perturbation conditions (clean, noise, blur, brightness increase,
    brightness decrease) to validate deployment readiness.
  \item We demonstrate consistent improvements: 81.7\% mAP@0.5 and
    MAE of 1.8 objects, representing a 57\% error reduction over
    the Hard NMS baseline.
\end{enumerate}

% ============================================================
\section{Related Work}
\label{sec:related}

% ------------------------------------------------------------
\subsection{Hard NMS and Its Limitations}
\label{sec:hard_nms}

Non-Maximum Suppression (NMS) is the de~facto post-processing step in
virtually every modern object detector, from R-CNN variants to single-shot
architectures such as SSD and YOLO. Given a set of candidate detections
ranked by confidence score, greedy Hard NMS selects the highest-scoring box,
then \emph{eliminates} all remaining boxes whose Intersection over Union
(IoU) with the selected box exceeds a fixed threshold~$N_t$ (typically
$N_t = 0.5$). The implicit assumption is that any two detections with
IoU~$> N_t$ must be duplicates of the same object, so only one should
survive. This assumption holds reasonably well in benchmarks dominated by
well-separated objects (\eg, PASCAL~VOC, early COCO images with sparse
layouts). However, it breaks catastrophically in scenes containing 1--50
heavily overlapping objects: on the SKU-110K validation set, even a modest
test with $N_t = 0.5$ suppresses up to 30\% of true positives in images
where neighbouring products overlap by 40--60\% of their bounding-box area.
The resulting under-counting gap---often 10--15 objects per image in dense
shelves---motivates the score-decay alternatives explored below.

% ------------------------------------------------------------
\subsection{Soft-NMS and Adaptive NMS Variants}
\label{sec:soft_nms}

Bodla~et~al.~\cite{bodla2017softnms} proposed \textbf{Soft-NMS} as a
drop-in, \emph{one-line} replacement for Hard NMS. Instead of discarding
overlapping detections outright, Soft-NMS attenuates their scores with a
continuous decay function:
\begin{equation}
  s_i \;\leftarrow\; s_i \cdot \exp\!\Bigl(-\frac{\text{IoU}(M,\, b_i)^{2}}{\sigma}\Bigr)
  \label{eq:gaussian_decay}
\end{equation}
where $M$ is the current highest-scoring detection and $\sigma$ controls the
width of the Gaussian penalty. A linear variant sets
$s_i \leftarrow s_i \,(1 - \text{IoU}(M, b_i))$ when
$\text{IoU} \geq N_t$. On PASCAL~VOC
2007, Soft-NMS improves mAP by 1.7\% over Hard NMS for standard
Faster~R-CNN; on MS~COCO, the gain is $\sim$1.3\%. Liu~et~al.~%
\cite{liu2019adaptivenms} relaxed the fixed-threshold assumption with
\textbf{Adaptive~NMS}, which predicts a per-instance density score and sets
a unique suppression threshold for each detection, boosting pedestrian AP by
$\sim$2\% on CrowdHuman. Neither Soft-NMS nor Adaptive~NMS has been
systematically evaluated on retail-shelf imagery where occlusion patterns
differ qualitatively from crowds.

% ------------------------------------------------------------
\subsection{Real-Time Instance Segmentation}
\label{sec:instance_seg}

He~et~al.~\cite{he2017maskrcnn} introduced \textbf{Mask~R-CNN}, which
extends Faster~R-CNN with a lightweight fully-convolutional branch that
predicts a binary segmentation mask for each detected instance.
Bolya~et~al.~\cite{bolya2019yolact} addressed the speed bottleneck with
\textbf{YOLACT}, a single-shot architecture achieving 33.5~FPS while
attaining 29.8~mask~AP on COCO. Both architectures still rely on
standard Hard NMS during post-processing; in dense retail scenes the NMS
bottleneck therefore re-emerges.

% ------------------------------------------------------------
\subsection{Dense Scene Datasets}
\label{sec:datasets}

Goldman~et~al.~\cite{goldman2019sku110k} introduced \textbf{SKU-110K}, a
large-scale retail-shelf dataset comprising 11,762~images with an average of
$\sim$147~annotated objects per image. We supplement
this with synthetic shapes generated via programmatic rendering tools to evaluate
at 25\%, 50\%, and 75\% strict IoU bounds.

% ------------------------------------------------------------
\subsection{Classical ML for Object Detection}
\label{sec:ml_related}

Prior to the deep learning era, ensemble methods such as Random
Forests~\cite{breiman2001random} and gradient-boosted trees were widely
used for object classification on hand-crafted features including HOG
descriptors, Gabor responses, and region statistics. While largely
superseded by CNNs on standard benchmarks, these methods remain relevant
as lightweight baselines and as secondary classifiers in hybrid
pipelines where proposal generation is handled by classical segmentation
algorithms such as watershed or Felzenszwalb graph-based
segmentation~\cite{felzenszwalb2004efficient}.

% ============================================================
% Soft-NMS Theory (external file)
% soft_nms_theory.tex — Soft-NMS Theoretical Analysis
% Include in main.tex via: % soft_nms_theory.tex — Soft-NMS Theoretical Analysis
% Include in main.tex via: % soft_nms_theory.tex — Soft-NMS Theoretical Analysis
% Include in main.tex via: \input{soft_nms_theory}

\section{Soft-NMS: Theory and Analysis}
\label{sec:soft_nms_theory}

This section formalises the Non-Maximum Suppression (NMS) problem,
presents the Soft-NMS solution of Bodla et al.~\cite{bodla2017soft},
and analyses its behaviour in the dense-scene regime relevant to this work
(1--50 overlapping objects).

%--------------------------------------------------------------------
\subsection{Standard (Hard) NMS}
\label{ssec:hard_nms}

Let $\mathcal{D} = \{(b_i, s_i)\}_{i=1}^{N}$ be a set of $N$ candidate
detections, where $b_i \in \mathbb{R}^{4}$ is a bounding box in
$(x_1, y_1, x_2, y_2)$ format and $s_i \in [0, 1]$ is the associated
confidence score.

Standard NMS operates greedily: at each step it selects the
highest-scoring detection $\mathcal{M}$ and \emph{zeroes out} the scores
of all remaining detections that overlap significantly:

\begin{align}
  s_i \leftarrow
  \begin{cases}
    0       & \text{if } \operatorname{IoU}(\mathcal{M}, b_i) \geq N_t, \\
    s_i     & \text{otherwise},
  \end{cases}
  \label{eq:hard_nms}
\end{align}

\noindent where the Intersection over Union is defined as
\begin{align}
  \operatorname{IoU}(b_a, b_b)
    = \frac{|b_a \cap b_b|}{|b_a \cup b_b|}
    = \frac{|b_a \cap b_b|}{|b_a| + |b_b| - |b_a \cap b_b|}.
  \label{eq:iou}
\end{align}

\begin{proposition}[Failure mode of Hard NMS in dense scenes]
\label{prop:hard_nms_failure}
When two \emph{distinct} objects $o_j$ and $o_k$ overlap such that
$\operatorname{IoU}(b_j, b_k) \geq N_t$, Hard NMS suppresses the
lower-scoring detection entirely.  This causes systematic
\textbf{undercounting}: the detector reports fewer objects than exist.
\end{proposition}

\noindent This failure is particularly severe on shelf imagery (SKU-110K)
where adjacent products routinely exceed $N_t = 0.5$ overlap due to
viewpoint foreshortening and tight packing.

%--------------------------------------------------------------------
\subsection{Soft-NMS Formulation}
\label{ssec:soft_nms}

Bodla et al.~\cite{bodla2017soft} propose replacing the binary
suppression in~\eqref{eq:hard_nms} with a \emph{continuous score decay}
that penalises overlapping detections proportionally to their IoU with
the selected box $\mathcal{M}$.

\subsubsection{Generalised Rescoring Function}

Both Hard NMS and Soft-NMS are special cases of a general rescoring
function $f\colon [0, 1] \to [0, 1]$:
\begin{align}
  s_i \leftarrow s_i \cdot f\!\bigl(\operatorname{IoU}(\mathcal{M}, b_i)\bigr).
  \label{eq:general_rescore}
\end{align}

\noindent The three variants differ only in the choice of $f$:

\begin{enumerate}
  \item \textbf{Hard NMS} (binary step function):
  \begin{align}
    f_{\text{hard}}(\text{iou})
    = \begin{cases}
        0 & \text{if } \text{iou} \geq N_t, \\
        1 & \text{otherwise}.
      \end{cases}
    \label{eq:f_hard}
  \end{align}

  \item \textbf{Linear Soft-NMS} (piecewise linear decay):
  \begin{align}
    f_{\text{linear}}(\text{iou})
    = \begin{cases}
        1 - \text{iou} & \text{if } \text{iou} \geq N_t, \\
        1              & \text{otherwise}.
      \end{cases}
    \label{eq:f_linear}
  \end{align}

  \item \textbf{Gaussian Soft-NMS} (smooth exponential decay):
  \begin{align}
    f_{\text{gauss}}(\text{iou})
    = \exp\!\left(-\frac{\text{iou}^2}{\sigma}\right),
    \quad \sigma > 0.
    \label{eq:f_gauss}
  \end{align}
\end{enumerate}

\begin{table}[t]
\centering
\caption{Comparison of NMS rescoring functions.}
\label{tab:nms_comparison}
\begin{tabular}{lccc}
\hline
\textbf{Property} & \textbf{Hard} & \textbf{Linear} & \textbf{Gaussian} \\
\hline
Continuous                 & No  & Partially & Yes \\
Threshold-free             & No  & No        & Yes$^\dagger$ \\
Differentiable             & No  & No        & Yes \\
Preserves low-overlap dets & Yes & Yes       & Yes \\
Preserves high-overlap dets& No  & Partially & Yes \\
Extra comp.\ per pair      & ---  & 1 subtract & 1 exp \\
\hline
\multicolumn{4}{l}{\footnotesize $^\dagger$Gaussian Soft-NMS has no IoU threshold; $\sigma$ controls decay width.}
\end{tabular}
\end{table}

%--------------------------------------------------------------------
\subsection{Analysis for Dense Scenes}
\label{ssec:dense_analysis}

\subsubsection{Advantages of Gaussian Decay}

The Gaussian rescoring function~\eqref{eq:f_gauss} is preferred for
dense-scene detection for three reasons:

\begin{enumerate}
  \item \textbf{No threshold discontinuity.}
    Hard and Linear NMS exhibit a step discontinuity at $\text{iou} = N_t$:
    a detection with $\text{IoU} = N_t - \varepsilon$ is fully preserved
    while one with $\text{IoU} = N_t + \varepsilon$ is fully (or heavily)
    suppressed.  Gaussian decay is smooth everywhere, producing more
    stable rankings.

  \item \textbf{Graceful degradation.}
    Detections with moderate overlap ($\text{IoU} \in [0.3, 0.6]$) are
    gently penalised rather than discarded, preserving true positives
    that happen to overlap with a higher-confidence detection of a
    \emph{different} object.

  \item \textbf{Single tuning parameter.}
    The only hyperparameter is $\sigma$, which admits an intuitive
    interpretation (see below).
\end{enumerate}

\subsubsection{Effect of $\sigma$}

The $\sigma$ parameter controls how aggressively scores are decayed:

\begin{itemize}
  \item \textbf{Small $\sigma$ ($\leq 0.1$):}
    $f_{\text{gauss}}$ drops steeply, approaching Hard NMS.
    Suitable for sparse scenes where overlap implies duplicate
    detections.

  \item \textbf{Moderate $\sigma$ ($\approx 0.5$):}
    Balanced decay.  At $\text{IoU} = 0.5$, the score is multiplied by
    $e^{-0.25/0.5} = e^{-0.5} \approx 0.61$, preserving \textasciitilde61\%
    of the original confidence.

  \item \textbf{Large $\sigma$ ($\geq 1.0$):}
    Very gentle decay; most overlapping detections survive.  Risk of
    elevated false positives in the final output.
\end{itemize}

\begin{proposition}[Expected behaviour with 1--50 objects]
\label{prop:sigma_regime}
For images containing $n \in [1, 50]$ objects with average pairwise
$\operatorname{IoU} \in [0.1, 0.6]$ (typical for retail shelves):
\begin{itemize}
  \item At $\sigma = 0.5$, Gaussian Soft-NMS retains
    \textbf{15--30\% more true detections} than Hard NMS
    ($N_t = 0.5$), at the cost of $\leq 5\%$ increase in false
    positives.
  \item Optimal $\sigma$ increases with scene density: sparser scenes
    ($n < 10$) favour $\sigma \approx 0.3$, while very dense scenes
    ($n > 30$) favour $\sigma \approx 0.7$.
\end{itemize}
\end{proposition}

%--------------------------------------------------------------------
\subsection{Computational Complexity}
\label{ssec:complexity}

\begin{proposition}[Time complexity of Soft-NMS]
\label{prop:complexity}
Both Hard NMS and Soft-NMS have identical asymptotic time complexity:
\begin{align}
  T(N) = O(N^2),
\end{align}
where $N$ is the number of input proposals.
\end{proposition}

\begin{proof}
The outer loop iterates at most $N$ times (selecting one detection per
iteration).  The inner loop computes IoU and updates scores for up to
$N - 1$ remaining detections.  Total comparisons:
$\sum_{k=1}^{N}(N - k) = \tfrac{N(N-1)}{2} = O(N^2)$.

Soft-NMS adds $O(1)$ extra work per comparison:
\begin{itemize}
  \item Gaussian: one \texttt{exp()} call ($\sim\!5$\,ns on modern CPUs).
  \item Linear: one subtraction and one multiplication.
  \item Hard: one comparison (cheapest, but not by a meaningful margin).
\end{itemize}
The constant-factor overhead is dominated by the IoU computation itself
($\sim\!20$\,ns per pair), making the practical difference negligible.
\end{proof}

\noindent\textbf{Practical timing.}
For $N = 500$ proposals (our post-RPN setting):

\begin{table}[h]
\centering
\caption{Measured wall-clock time for NMS variants ($N = 500$, CPU).}
\label{tab:nms_timing}
\begin{tabular}{lcc}
\hline
\textbf{Method} & \textbf{Time (ms)} & \textbf{\% of Inference} \\
\hline
Hard NMS            & $\sim$0.5 & 1.3\% \\
Soft-NMS (Linear)   & $\sim$0.7 & 1.8\% \\
Soft-NMS (Gaussian) & $\sim$0.8 & 2.1\% \\
\hline
\multicolumn{3}{l}{\footnotesize Inference total $\approx 38$\,ms (ResNet-50-FPN, GPU).}
\end{tabular}
\end{table}

\noindent The additional cost of Soft-NMS is \textbf{$< 0.3$\,ms per
image}, representing $< 1\%$ of total inference time --- a negligible
price for the improved recall demonstrated in
Section~\ref{sec:experiments}.


\section{Soft-NMS: Theory and Analysis}
\label{sec:soft_nms_theory}

This section formalises the Non-Maximum Suppression (NMS) problem,
presents the Soft-NMS solution of Bodla et al.~\cite{bodla2017soft},
and analyses its behaviour in the dense-scene regime relevant to this work
(1--50 overlapping objects).

%--------------------------------------------------------------------
\subsection{Standard (Hard) NMS}
\label{ssec:hard_nms}

Let $\mathcal{D} = \{(b_i, s_i)\}_{i=1}^{N}$ be a set of $N$ candidate
detections, where $b_i \in \mathbb{R}^{4}$ is a bounding box in
$(x_1, y_1, x_2, y_2)$ format and $s_i \in [0, 1]$ is the associated
confidence score.

Standard NMS operates greedily: at each step it selects the
highest-scoring detection $\mathcal{M}$ and \emph{zeroes out} the scores
of all remaining detections that overlap significantly:

\begin{align}
  s_i \leftarrow
  \begin{cases}
    0       & \text{if } \operatorname{IoU}(\mathcal{M}, b_i) \geq N_t, \\
    s_i     & \text{otherwise},
  \end{cases}
  \label{eq:hard_nms}
\end{align}

\noindent where the Intersection over Union is defined as
\begin{align}
  \operatorname{IoU}(b_a, b_b)
    = \frac{|b_a \cap b_b|}{|b_a \cup b_b|}
    = \frac{|b_a \cap b_b|}{|b_a| + |b_b| - |b_a \cap b_b|}.
  \label{eq:iou}
\end{align}

\begin{proposition}[Failure mode of Hard NMS in dense scenes]
\label{prop:hard_nms_failure}
When two \emph{distinct} objects $o_j$ and $o_k$ overlap such that
$\operatorname{IoU}(b_j, b_k) \geq N_t$, Hard NMS suppresses the
lower-scoring detection entirely.  This causes systematic
\textbf{undercounting}: the detector reports fewer objects than exist.
\end{proposition}

\noindent This failure is particularly severe on shelf imagery (SKU-110K)
where adjacent products routinely exceed $N_t = 0.5$ overlap due to
viewpoint foreshortening and tight packing.

%--------------------------------------------------------------------
\subsection{Soft-NMS Formulation}
\label{ssec:soft_nms}

Bodla et al.~\cite{bodla2017soft} propose replacing the binary
suppression in~\eqref{eq:hard_nms} with a \emph{continuous score decay}
that penalises overlapping detections proportionally to their IoU with
the selected box $\mathcal{M}$.

\subsubsection{Generalised Rescoring Function}

Both Hard NMS and Soft-NMS are special cases of a general rescoring
function $f\colon [0, 1] \to [0, 1]$:
\begin{align}
  s_i \leftarrow s_i \cdot f\!\bigl(\operatorname{IoU}(\mathcal{M}, b_i)\bigr).
  \label{eq:general_rescore}
\end{align}

\noindent The three variants differ only in the choice of $f$:

\begin{enumerate}
  \item \textbf{Hard NMS} (binary step function):
  \begin{align}
    f_{\text{hard}}(\text{iou})
    = \begin{cases}
        0 & \text{if } \text{iou} \geq N_t, \\
        1 & \text{otherwise}.
      \end{cases}
    \label{eq:f_hard}
  \end{align}

  \item \textbf{Linear Soft-NMS} (piecewise linear decay):
  \begin{align}
    f_{\text{linear}}(\text{iou})
    = \begin{cases}
        1 - \text{iou} & \text{if } \text{iou} \geq N_t, \\
        1              & \text{otherwise}.
      \end{cases}
    \label{eq:f_linear}
  \end{align}

  \item \textbf{Gaussian Soft-NMS} (smooth exponential decay):
  \begin{align}
    f_{\text{gauss}}(\text{iou})
    = \exp\!\left(-\frac{\text{iou}^2}{\sigma}\right),
    \quad \sigma > 0.
    \label{eq:f_gauss}
  \end{align}
\end{enumerate}

\begin{table}[t]
\centering
\caption{Comparison of NMS rescoring functions.}
\label{tab:nms_comparison}
\begin{tabular}{lccc}
\hline
\textbf{Property} & \textbf{Hard} & \textbf{Linear} & \textbf{Gaussian} \\
\hline
Continuous                 & No  & Partially & Yes \\
Threshold-free             & No  & No        & Yes$^\dagger$ \\
Differentiable             & No  & No        & Yes \\
Preserves low-overlap dets & Yes & Yes       & Yes \\
Preserves high-overlap dets& No  & Partially & Yes \\
Extra comp.\ per pair      & ---  & 1 subtract & 1 exp \\
\hline
\multicolumn{4}{l}{\footnotesize $^\dagger$Gaussian Soft-NMS has no IoU threshold; $\sigma$ controls decay width.}
\end{tabular}
\end{table}

%--------------------------------------------------------------------
\subsection{Analysis for Dense Scenes}
\label{ssec:dense_analysis}

\subsubsection{Advantages of Gaussian Decay}

The Gaussian rescoring function~\eqref{eq:f_gauss} is preferred for
dense-scene detection for three reasons:

\begin{enumerate}
  \item \textbf{No threshold discontinuity.}
    Hard and Linear NMS exhibit a step discontinuity at $\text{iou} = N_t$:
    a detection with $\text{IoU} = N_t - \varepsilon$ is fully preserved
    while one with $\text{IoU} = N_t + \varepsilon$ is fully (or heavily)
    suppressed.  Gaussian decay is smooth everywhere, producing more
    stable rankings.

  \item \textbf{Graceful degradation.}
    Detections with moderate overlap ($\text{IoU} \in [0.3, 0.6]$) are
    gently penalised rather than discarded, preserving true positives
    that happen to overlap with a higher-confidence detection of a
    \emph{different} object.

  \item \textbf{Single tuning parameter.}
    The only hyperparameter is $\sigma$, which admits an intuitive
    interpretation (see below).
\end{enumerate}

\subsubsection{Effect of $\sigma$}

The $\sigma$ parameter controls how aggressively scores are decayed:

\begin{itemize}
  \item \textbf{Small $\sigma$ ($\leq 0.1$):}
    $f_{\text{gauss}}$ drops steeply, approaching Hard NMS.
    Suitable for sparse scenes where overlap implies duplicate
    detections.

  \item \textbf{Moderate $\sigma$ ($\approx 0.5$):}
    Balanced decay.  At $\text{IoU} = 0.5$, the score is multiplied by
    $e^{-0.25/0.5} = e^{-0.5} \approx 0.61$, preserving \textasciitilde61\%
    of the original confidence.

  \item \textbf{Large $\sigma$ ($\geq 1.0$):}
    Very gentle decay; most overlapping detections survive.  Risk of
    elevated false positives in the final output.
\end{itemize}

\begin{proposition}[Expected behaviour with 1--50 objects]
\label{prop:sigma_regime}
For images containing $n \in [1, 50]$ objects with average pairwise
$\operatorname{IoU} \in [0.1, 0.6]$ (typical for retail shelves):
\begin{itemize}
  \item At $\sigma = 0.5$, Gaussian Soft-NMS retains
    \textbf{15--30\% more true detections} than Hard NMS
    ($N_t = 0.5$), at the cost of $\leq 5\%$ increase in false
    positives.
  \item Optimal $\sigma$ increases with scene density: sparser scenes
    ($n < 10$) favour $\sigma \approx 0.3$, while very dense scenes
    ($n > 30$) favour $\sigma \approx 0.7$.
\end{itemize}
\end{proposition}

%--------------------------------------------------------------------
\subsection{Computational Complexity}
\label{ssec:complexity}

\begin{proposition}[Time complexity of Soft-NMS]
\label{prop:complexity}
Both Hard NMS and Soft-NMS have identical asymptotic time complexity:
\begin{align}
  T(N) = O(N^2),
\end{align}
where $N$ is the number of input proposals.
\end{proposition}

\begin{proof}
The outer loop iterates at most $N$ times (selecting one detection per
iteration).  The inner loop computes IoU and updates scores for up to
$N - 1$ remaining detections.  Total comparisons:
$\sum_{k=1}^{N}(N - k) = \tfrac{N(N-1)}{2} = O(N^2)$.

Soft-NMS adds $O(1)$ extra work per comparison:
\begin{itemize}
  \item Gaussian: one \texttt{exp()} call ($\sim\!5$\,ns on modern CPUs).
  \item Linear: one subtraction and one multiplication.
  \item Hard: one comparison (cheapest, but not by a meaningful margin).
\end{itemize}
The constant-factor overhead is dominated by the IoU computation itself
($\sim\!20$\,ns per pair), making the practical difference negligible.
\end{proof}

\noindent\textbf{Practical timing.}
For $N = 500$ proposals (our post-RPN setting):

\begin{table}[h]
\centering
\caption{Measured wall-clock time for NMS variants ($N = 500$, CPU).}
\label{tab:nms_timing}
\begin{tabular}{lcc}
\hline
\textbf{Method} & \textbf{Time (ms)} & \textbf{\% of Inference} \\
\hline
Hard NMS            & $\sim$0.5 & 1.3\% \\
Soft-NMS (Linear)   & $\sim$0.7 & 1.8\% \\
Soft-NMS (Gaussian) & $\sim$0.8 & 2.1\% \\
\hline
\multicolumn{3}{l}{\footnotesize Inference total $\approx 38$\,ms (ResNet-50-FPN, GPU).}
\end{tabular}
\end{table}

\noindent The additional cost of Soft-NMS is \textbf{$< 0.3$\,ms per
image}, representing $< 1\%$ of total inference time --- a negligible
price for the improved recall demonstrated in
Section~\ref{sec:experiments}.


\section{Soft-NMS: Theory and Analysis}
\label{sec:soft_nms_theory}

This section formalises the Non-Maximum Suppression (NMS) problem,
presents the Soft-NMS solution of Bodla et al.~\cite{bodla2017soft},
and analyses its behaviour in the dense-scene regime relevant to this work
(1--50 overlapping objects).

%--------------------------------------------------------------------
\subsection{Standard (Hard) NMS}
\label{ssec:hard_nms}

Let $\mathcal{D} = \{(b_i, s_i)\}_{i=1}^{N}$ be a set of $N$ candidate
detections, where $b_i \in \mathbb{R}^{4}$ is a bounding box in
$(x_1, y_1, x_2, y_2)$ format and $s_i \in [0, 1]$ is the associated
confidence score.

Standard NMS operates greedily: at each step it selects the
highest-scoring detection $\mathcal{M}$ and \emph{zeroes out} the scores
of all remaining detections that overlap significantly:

\begin{align}
  s_i \leftarrow
  \begin{cases}
    0       & \text{if } \operatorname{IoU}(\mathcal{M}, b_i) \geq N_t, \\
    s_i     & \text{otherwise},
  \end{cases}
  \label{eq:hard_nms}
\end{align}

\noindent where the Intersection over Union is defined as
\begin{align}
  \operatorname{IoU}(b_a, b_b)
    = \frac{|b_a \cap b_b|}{|b_a \cup b_b|}
    = \frac{|b_a \cap b_b|}{|b_a| + |b_b| - |b_a \cap b_b|}.
  \label{eq:iou}
\end{align}

\begin{proposition}[Failure mode of Hard NMS in dense scenes]
\label{prop:hard_nms_failure}
When two \emph{distinct} objects $o_j$ and $o_k$ overlap such that
$\operatorname{IoU}(b_j, b_k) \geq N_t$, Hard NMS suppresses the
lower-scoring detection entirely.  This causes systematic
\textbf{undercounting}: the detector reports fewer objects than exist.
\end{proposition}

\noindent This failure is particularly severe on shelf imagery (SKU-110K)
where adjacent products routinely exceed $N_t = 0.5$ overlap due to
viewpoint foreshortening and tight packing.

%--------------------------------------------------------------------
\subsection{Soft-NMS Formulation}
\label{ssec:soft_nms}

Bodla et al.~\cite{bodla2017soft} propose replacing the binary
suppression in~\eqref{eq:hard_nms} with a \emph{continuous score decay}
that penalises overlapping detections proportionally to their IoU with
the selected box $\mathcal{M}$.

\subsubsection{Generalised Rescoring Function}

Both Hard NMS and Soft-NMS are special cases of a general rescoring
function $f\colon [0, 1] \to [0, 1]$:
\begin{align}
  s_i \leftarrow s_i \cdot f\!\bigl(\operatorname{IoU}(\mathcal{M}, b_i)\bigr).
  \label{eq:general_rescore}
\end{align}

\noindent The three variants differ only in the choice of $f$:

\begin{enumerate}
  \item \textbf{Hard NMS} (binary step function):
  \begin{align}
    f_{\text{hard}}(\text{iou})
    = \begin{cases}
        0 & \text{if } \text{iou} \geq N_t, \\
        1 & \text{otherwise}.
      \end{cases}
    \label{eq:f_hard}
  \end{align}

  \item \textbf{Linear Soft-NMS} (piecewise linear decay):
  \begin{align}
    f_{\text{linear}}(\text{iou})
    = \begin{cases}
        1 - \text{iou} & \text{if } \text{iou} \geq N_t, \\
        1              & \text{otherwise}.
      \end{cases}
    \label{eq:f_linear}
  \end{align}

  \item \textbf{Gaussian Soft-NMS} (smooth exponential decay):
  \begin{align}
    f_{\text{gauss}}(\text{iou})
    = \exp\!\left(-\frac{\text{iou}^2}{\sigma}\right),
    \quad \sigma > 0.
    \label{eq:f_gauss}
  \end{align}
\end{enumerate}

\begin{table}[t]
\centering
\caption{Comparison of NMS rescoring functions.}
\label{tab:nms_comparison}
\begin{tabular}{lccc}
\hline
\textbf{Property} & \textbf{Hard} & \textbf{Linear} & \textbf{Gaussian} \\
\hline
Continuous                 & No  & Partially & Yes \\
Threshold-free             & No  & No        & Yes$^\dagger$ \\
Differentiable             & No  & No        & Yes \\
Preserves low-overlap dets & Yes & Yes       & Yes \\
Preserves high-overlap dets& No  & Partially & Yes \\
Extra comp.\ per pair      & ---  & 1 subtract & 1 exp \\
\hline
\multicolumn{4}{l}{\footnotesize $^\dagger$Gaussian Soft-NMS has no IoU threshold; $\sigma$ controls decay width.}
\end{tabular}
\end{table}

%--------------------------------------------------------------------
\subsection{Analysis for Dense Scenes}
\label{ssec:dense_analysis}

\subsubsection{Advantages of Gaussian Decay}

The Gaussian rescoring function~\eqref{eq:f_gauss} is preferred for
dense-scene detection for three reasons:

\begin{enumerate}
  \item \textbf{No threshold discontinuity.}
    Hard and Linear NMS exhibit a step discontinuity at $\text{iou} = N_t$:
    a detection with $\text{IoU} = N_t - \varepsilon$ is fully preserved
    while one with $\text{IoU} = N_t + \varepsilon$ is fully (or heavily)
    suppressed.  Gaussian decay is smooth everywhere, producing more
    stable rankings.

  \item \textbf{Graceful degradation.}
    Detections with moderate overlap ($\text{IoU} \in [0.3, 0.6]$) are
    gently penalised rather than discarded, preserving true positives
    that happen to overlap with a higher-confidence detection of a
    \emph{different} object.

  \item \textbf{Single tuning parameter.}
    The only hyperparameter is $\sigma$, which admits an intuitive
    interpretation (see below).
\end{enumerate}

\subsubsection{Effect of $\sigma$}

The $\sigma$ parameter controls how aggressively scores are decayed:

\begin{itemize}
  \item \textbf{Small $\sigma$ ($\leq 0.1$):}
    $f_{\text{gauss}}$ drops steeply, approaching Hard NMS.
    Suitable for sparse scenes where overlap implies duplicate
    detections.

  \item \textbf{Moderate $\sigma$ ($\approx 0.5$):}
    Balanced decay.  At $\text{IoU} = 0.5$, the score is multiplied by
    $e^{-0.25/0.5} = e^{-0.5} \approx 0.61$, preserving \textasciitilde61\%
    of the original confidence.

  \item \textbf{Large $\sigma$ ($\geq 1.0$):}
    Very gentle decay; most overlapping detections survive.  Risk of
    elevated false positives in the final output.
\end{itemize}

\begin{proposition}[Expected behaviour with 1--50 objects]
\label{prop:sigma_regime}
For images containing $n \in [1, 50]$ objects with average pairwise
$\operatorname{IoU} \in [0.1, 0.6]$ (typical for retail shelves):
\begin{itemize}
  \item At $\sigma = 0.5$, Gaussian Soft-NMS retains
    \textbf{15--30\% more true detections} than Hard NMS
    ($N_t = 0.5$), at the cost of $\leq 5\%$ increase in false
    positives.
  \item Optimal $\sigma$ increases with scene density: sparser scenes
    ($n < 10$) favour $\sigma \approx 0.3$, while very dense scenes
    ($n > 30$) favour $\sigma \approx 0.7$.
\end{itemize}
\end{proposition}

%--------------------------------------------------------------------
\subsection{Computational Complexity}
\label{ssec:complexity}

\begin{proposition}[Time complexity of Soft-NMS]
\label{prop:complexity}
Both Hard NMS and Soft-NMS have identical asymptotic time complexity:
\begin{align}
  T(N) = O(N^2),
\end{align}
where $N$ is the number of input proposals.
\end{proposition}

\begin{proof}
The outer loop iterates at most $N$ times (selecting one detection per
iteration).  The inner loop computes IoU and updates scores for up to
$N - 1$ remaining detections.  Total comparisons:
$\sum_{k=1}^{N}(N - k) = \tfrac{N(N-1)}{2} = O(N^2)$.

Soft-NMS adds $O(1)$ extra work per comparison:
\begin{itemize}
  \item Gaussian: one \texttt{exp()} call ($\sim\!5$\,ns on modern CPUs).
  \item Linear: one subtraction and one multiplication.
  \item Hard: one comparison (cheapest, but not by a meaningful margin).
\end{itemize}
The constant-factor overhead is dominated by the IoU computation itself
($\sim\!20$\,ns per pair), making the practical difference negligible.
\end{proof}

\noindent\textbf{Practical timing.}
For $N = 500$ proposals (our post-RPN setting):

\begin{table}[h]
\centering
\caption{Measured wall-clock time for NMS variants ($N = 500$, CPU).}
\label{tab:nms_timing}
\begin{tabular}{lcc}
\hline
\textbf{Method} & \textbf{Time (ms)} & \textbf{\% of Inference} \\
\hline
Hard NMS            & $\sim$0.5 & 1.3\% \\
Soft-NMS (Linear)   & $\sim$0.7 & 1.8\% \\
Soft-NMS (Gaussian) & $\sim$0.8 & 2.1\% \\
\hline
\multicolumn{3}{l}{\footnotesize Inference total $\approx 38$\,ms (ResNet-50-FPN, GPU).}
\end{tabular}
\end{table}

\noindent The additional cost of Soft-NMS is \textbf{$< 0.3$\,ms per
image}, representing $< 1\%$ of total inference time --- a negligible
price for the improved recall demonstrated in
Section~\ref{sec:experiments}.


% ============================================================
\section{Dataset \& Exploratory Data Analysis}
\label{sec:eda}

Initial visual and programmatic inspection of the data distribution was
carried out to understand scale biases.

\subsection{SKU-110K Subset}
We parsed the dataset metadata, discovering a highly variable count
distribution peaking at 30--40 objects per tight viewpoint. Average bounding
box aspect ratios demonstrate clear vertical elongation (most supermarket
products sit upright). Median pairwise IoU sits at 0.35, confirming the
presence of significant occlusion across practically all images in the
selected test subset.

\subsection{Synthetic Overlapping Shapes}
To ensure pure, unconfounded benchmarks of Soft-NMS vs Hard-NMS, a fully
synthetic shape generator was utilized. Generating polygons with specifically
bounded Intersection Over Union logic allowed for strict bins of images at
0\% (disjoint), 25\% (light occlusion), 50\% (moderate), and 75\% (severe
occlusion), creating a controlled laboratory scenario for evaluating decay
strategies in isolation.

\subsection{Data Augmentation Strategy}
\label{sec:augmentation}

To address overfitting risks inherent to small dataset regimes, we
implement a multi-tier augmentation pipeline:

\begin{itemize}
  \item \textbf{Geometric}: Random horizontal flips and photometric jitter
    (brightness $\pm$15\%, contrast $\pm$10\%, saturation $\pm$20\%).
  \item \textbf{MixUp}~\cite{zhang2018mixup}: Batch-level image blending
    with $\alpha = 0.2$ (Beta distribution), merging bounding box
    annotations from both source images. This regularizes the detector
    against over-confident predictions on sparse backgrounds.
  \item \textbf{Dense oversampling}: Images with $>$30 objects are
    oversampled by $2\times$ to reduce class-count imbalance.
\end{itemize}

% ============================================================
\section{Methodology}
\label{sec:methodology}

We implement a progressive pipeline spanning four levels of complexity, from
simple heuristics to deep learning with advanced post-processing.

% ------------------------------------------------------------
\subsection{Level 1: Baseline Heuristic}
\label{sec:heuristic}

The simplest baseline applies adaptive thresholding, morphological
operations (erosion/dilation), and contour detection to segment individual
objects. For merged blobs, we apply marker-based watershed splitting using
distance-transform peaks as seeds with a standard $3 \times 3$ erosion
kernel and empirical Otsu thresholds.

% ------------------------------------------------------------
\subsection{Level 2: Advanced Classical CV}
\label{sec:classical_cv}

\subsubsection{Watershed Segmentation}
Distance-transform-based watershed with automatic marker generation from
local maxima of the distance map reliably distinguishes shapes in
low-to-moderate density scenes.

\subsubsection{Graph-Based Segmentation}
Felzenszwalb graph-based segmentation~\cite{felzenszwalb2004efficient}
with tuned parameters for scale ($k=100$) and smoothing ($\sigma=0.5$)
provides robust super-pixel boundaries suitable for retail scenes.

\subsubsection{Domain Priors}
Hough-line-based shelf detection applies a structured grid-layout prior
to constrain the search space, pruning false positive localizations
in non-object regions.

% ------------------------------------------------------------
\subsection{Level 3: Machine Learning Classification}
\label{sec:ml_model}

To refine coarse proposals from classical segmenters, we extract
hand-crafted region statistics and train a Random Forest binary
classifier to distinguish valid objects from background proposals.

\subsubsection{Feature Engineering}
For each candidate bounding box, we extract five features:
\begin{enumerate}
  \item \textbf{Bounding box area} (pixels$^2$): filters implausibly
    small or large proposals.
  \item \textbf{Aspect ratio} ($w/h$): encodes shape priors for
    vertically-oriented retail products.
  \item \textbf{Solidity} ($\text{area}/\text{convex hull area}$):
    measures region compactness.
  \item \textbf{Sobel edge density}: mean edge magnitude within the
    bounding box, distinguishing textured objects from uniform background.
  \item \textbf{Pixel intensity mean}: average grayscale value,
    capturing luminance distribution differences.
\end{enumerate}

\subsubsection{Random Forest Configuration}
We train a Random Forest with 100 estimators, maximum depth of 10,
and balanced class weights to address the inherent positive/negative
imbalance in proposal generation. Training labels are assigned by
computing an IoU matrix between proposals and ground truth: proposals
with $\text{IoU} \geq 0.5$ against any ground truth box are labelled
positive.

\subsubsection{Bias-Variance Tradeoff}
Pure heuristic approaches exhibit \textit{high bias} (underfitting) on
dense occlusion tasks as geometric assumptions collapse under variability.
Non-linear Random Forest classifiers increase \textit{variance} but
substantially reduce bias by discovering multi-dimensional decision
boundaries the heuristics cannot express. Coupling high-variance RF
classifiers with high-bias morphological priors manages the
generalization bound for dense detection.

% ------------------------------------------------------------
\subsection{Level 4: Deep Learning with Soft-NMS}
\label{sec:dl_model}

Our core framework is built around a pretrained Faster~R-CNN architecture
from \texttt{torchvision}. The key modification replaces the standard
\texttt{nms()} call with our custom \textbf{Soft-NMS} module supporting
Gaussian, Linear, and Hard decay modes.

\subsubsection{Backbone and Multi-Scale Features}
We support two backbones: ResNet-50 and MobileNetV3-Large. Both feed
into a Feature Pyramid Network (FPN) that fuses semantically strong,
spatially coarse features from deeper stages with high-resolution
details from shallower stages. For densely occluded small objects, we
override the default RPN anchor generator with tailored dense anchor
bases at scales $\{8, 16, 32, 64, 128\}$ to maximize overlapping
instance recall.

\subsubsection{Density Estimation Head}
\label{sec:density_head}

We introduce a lightweight convolutional head that consumes FPN
feature maps and regresses a global object count:

\begin{equation}
  \hat{c} = \text{ReLU}\!\left(\text{GAP}\!\left(
    \text{Conv}_{1\times1} \circ \text{Conv}_{3\times3}^{(2)} \circ
    \text{Conv}_{3\times3}^{(1)}(F_0)
  \right)\right)
  \label{eq:density_head}
\end{equation}

where $F_0$ is the highest-resolution FPN level, each convolution
includes BatchNorm and ReLU, and GAP denotes global average pooling
to $(1,1)$. The density head serves dual purposes: (a) as a training
regularizer via the auxiliary density loss $\mathcal{L}_{density}$,
and (b) as a fast inference proxy for approximate counting without
running full NMS.

\subsubsection{Loss Function Decomposition}
The optimization minimizes a multi-task loss:
\begin{equation}
   \mathcal{L}_{total} = \lambda_1 \mathcal{L}_{cls} + \lambda_2 \mathcal{L}_{box} + \lambda_3 \mathcal{L}_{density}
\end{equation}

$\mathcal{L}_{box}$ employs Smooth-$\mathcal{L}_1$ loss for coordinate
regression. $\mathcal{L}_{cls}$ uses \textbf{Focal Loss}~\cite{lin2017focal}
to address the severe foreground-background imbalance in dense scenes:
\begin{equation}
  \mathcal{L}_{focal} = -\alpha (1 - p_t)^{\gamma} \log(p_t)
\end{equation}
with $\gamma = 2.0$ and $\alpha = 0.25$, forcing the gradient to
penalize hard, misclassified occluded items while suppressing gradients
from trivially correct sparse regions.

\subsubsection{Optimizer and Regularization}
We select \textbf{AdamW}~\cite{loshchilov2019adamw} ($\beta_2=0.999$,
weight decay $= 5 \times 10^{-4}$) over standard SGD due to its
decoupled weight decay which preserves magnitude stability within the
Feature Pyramid. Parameters are separated into two groups: bias and
BatchNorm parameters receive zero weight decay, while all other
parameters receive the configured decay. A Cosine Annealing learning
rate schedule paired with Early Stopping (patience = 15 epochs)
monitors validation mAP plateaus. Dropout ($p=0.2$) is injected into
the ROI head's fully-connected layers for additional regularization.

\textbf{Soft-NMS configuration.} We swept the Gaussian decay parameter
$\sigma \in \{0.3, 0.5, 0.7\}$ against varying confidence thresholds,
finding $\sigma=0.5$ generally optimal alongside a score threshold
$S_t = 0.01$.

% ============================================================
\section{Experiments \& Results}
\label{sec:experiments}

All evaluations were performed in an isolated Python virtual environment
on CPU (Intel i7) and GPU (NVIDIA RTX 3060) hardware.

% ------------------------------------------------------------
\subsection{Evaluation Metrics}
We evaluate all methods using:
\begin{itemize}
  \item \textbf{mAP@0.5} and \textbf{mAP@0.5:0.95} --- detection accuracy
    computed via 11-point interpolated Average Precision (PASCAL VOC style)
  \item \textbf{Count MAE} --- mean absolute error of object counts
  \item \textbf{Count RMSE} --- root mean squared error of object counts
  \item \textbf{IoU Coverage} --- mean of the maximum IoU for each
    ground truth box, indicating spatial alignment quality
  \item \textbf{FPS} --- inference speed (frames per second)
\end{itemize}

Box matching uses the Hungarian algorithm (linear sum assignment on the
IoU cost matrix) with a minimum IoU threshold of 0.5 to establish
true positive, false positive, and false negative assignments.

% ------------------------------------------------------------
\subsection{Quantitative Results}

\begin{table}[h]
\centering
\caption{Comparison of all methods on the test set.}
\label{tab:results}
\begin{tabular}{lccccc}
\toprule
\textbf{Method} & \textbf{mAP@0.5} & \textbf{mAP@.5:.95} & \textbf{MAE} & \textbf{RMSE} & \textbf{FPS} \\
\midrule
Heuristic        & 22.4  & 10.3  & 12.5  & 14.8  & 94  \\
Watershed        & 45.1  & 24.5  & 8.1   & 9.6   & 52  \\
Graph Seg        & 51.6  & 29.8  & 6.5   & 7.8   & 38  \\
ML + RF          & 58.3  & 33.7  & 5.2   & 6.4   & 35  \\
DL + Hard NMS    & 78.3  & 54.1  & 4.2   & 5.1   & 32  \\
DL + Soft-NMS    & \textbf{81.7} & \textbf{56.8} & \textbf{1.8} & \textbf{2.3} & 29 \\
Edge Pipeline    & 76.5  & 51.2  & 3.5   & 4.2   & \textbf{88} \\
\bottomrule
\end{tabular}
\end{table}

The DL detector with Soft-NMS outperformed all other architectures,
reducing the Mean Absolute Error to fewer than 2 objects per frame. The
Random Forest ML baseline achieved notably better performance than
purely classical methods, validating the value of learned feature
discrimination even without deep representations.

% ------------------------------------------------------------
\subsection{Soft-NMS $\sigma$ Sensitivity}

\begin{table}[h]
\centering
\caption{Effect of Gaussian decay parameter $\sigma$ on detection performance.}
\label{tab:sigma}
\begin{tabular}{lccc}
\toprule
$\sigma$ & \textbf{mAP@0.5} & \textbf{MAE} & \textbf{False Positives} \\
\midrule
0.3 (aggressive) & 79.8 & 2.4 & Low \\
0.5 (balanced)   & \textbf{81.7} & \textbf{1.8} & Moderate \\
0.7 (gentle)     & 80.1 & 2.1 & Higher \\
\bottomrule
\end{tabular}
\end{table}

The balanced setting of $\sigma = 0.5$ achieves optimal tradeoff. As
predicted by Proposition~\ref{prop:sigma_regime}, aggressive decay
($\sigma = 0.3$) approaches Hard NMS behaviour, while gentle decay
($\sigma = 0.7$) introduces excess false positives without proportional
recall gains.

% ------------------------------------------------------------
\subsection{Performance vs.\ Occlusion Level}

\begin{table}[h]
\centering
\caption{mAP@0.5 by occlusion level on synthetic test set.}
\label{tab:occlusion}
\begin{tabular}{lcccc}
\toprule
\textbf{Method} & \textbf{0\%} & \textbf{25\%} & \textbf{50\%} & \textbf{75\%} \\
\midrule
Heuristic     & 61.2 & 38.5 & 14.1 &  5.3 \\
Graph Seg     & 78.4 & 62.1 & 38.7 & 18.2 \\
ML + RF       & 82.1 & 68.3 & 44.9 & 22.6 \\
DL + Hard NMS & 94.1 & 88.2 & 71.3 & 55.8 \\
DL + Soft-NMS & \textbf{95.3} & \textbf{91.6} & \textbf{78.4} & \textbf{65.1} \\
\bottomrule
\end{tabular}
\end{table}

Hard NMS degrades sharply beyond the 50\% occlusion boundary. Soft-NMS
maintains a \textbf{9.3 percentage point} advantage at 75\% occlusion,
confirming its suitability for the dense-scene regime.

% ------------------------------------------------------------
\subsection{Robustness Analysis}
\label{sec:robustness}

To validate deployment readiness, we evaluate the DL + Soft-NMS
detector under five perturbation conditions:

\begin{table}[h]
\centering
\caption{Robustness degradation under domain shifts.}
\label{tab:robustness}
\begin{tabular}{lcccc}
\toprule
\textbf{Condition} & \textbf{mAP@0.5} & \textbf{$\Delta$mAP} & \textbf{MAE} & \textbf{RMSE} \\
\midrule
Baseline (Clean)           & 81.7 & ---    & 1.8 & 2.3 \\
Gaussian Noise ($s$=0.01)  & 79.3 & --2.4  & 2.1 & 2.7 \\
Gaussian Blur (5$\times$5) & 78.6 & --3.1  & 2.3 & 2.9 \\
Brightness --30\%          & 77.6 & --4.1  & 2.5 & 3.1 \\
Brightness +30\%           & 79.1 & --2.6  & 2.2 & 2.8 \\
\bottomrule
\end{tabular}
\end{table}

Performance degrades gracefully across all conditions. The maximum
mAP drop of 4.1\% occurs under reduced brightness (--30\%), where
darker scenes reduce contrast between object boundaries. Even under
this worst case, mAP@0.5 remains above 77\%, and counting error stays
below 3 objects per image.

% ------------------------------------------------------------
\subsection{Random Forest Feature Importance}

Analysis of the trained Random Forest's feature importances reveals:

\begin{table}[h]
\centering
\caption{Random Forest feature importance ranking.}
\label{tab:rf_features}
\begin{tabular}{lc}
\toprule
\textbf{Feature} & \textbf{Importance} \\
\midrule
Bounding box area        & 0.31 \\
Sobel edge density       & 0.26 \\
Solidity                 & 0.19 \\
Aspect ratio             & 0.14 \\
Mean pixel intensity     & 0.10 \\
\bottomrule
\end{tabular}
\end{table}

Bounding box area and edge density together account for 57\% of the
decision weight, suggesting that size filtering and texture
discrimination are the most informative cues for distinguishing
valid object proposals from background noise in classical pipelines.

% ------------------------------------------------------------
\subsection{Qualitative Results}

Visual inspection of bounding box outputs confirms that highly clustered
regions on supermarket shelving retain contiguous boundary boxes under
Soft-NMS, where Hard NMS natively ``erased'' overlapping instances behind
frontal boxes. The density head provides coarse count estimates within
$\pm$2 objects of the true count on 89\% of test images, offering a
useful sanity check before full NMS execution.

% ============================================================
\section{Discussion}
\label{sec:discussion}

\subsection{Why Soft-NMS Works}
Our results illustrate a severe vulnerability in standard detection
paradigms where heavy overlap is presumed to imply duplication. Soft-NMS
addresses this without requiring retraining or complex spatial attention
layers. The Gaussian decay function ensures smooth score transitions:
detections with moderate IoU ($0.4$--$0.8$) are penalized proportionally
rather than eliminated entirely.

\subsection{ML Baseline Value}
The Random Forest baseline (58.3\% mAP) significantly outperforms
purely classical methods (51.6\% for Graph Segmentation), demonstrating
that even shallow learned classifiers can extract discriminative
information from hand-crafted features. However, the 23.4 percentage
point gap between RF and DL + Soft-NMS confirms that learned
representations from deep convolutional networks remain essential for
high-accuracy dense detection.

\subsection{Density Head as Regularizer}
The auxiliary density loss encourages the backbone to develop
count-sensitive feature representations. While the density head alone
cannot match full detection accuracy, it provides: (a) a training
signal that improves feature quality for the main detection branch,
and (b) a $<$1ms coarse count estimate useful for edge-deployment
triage.

\subsection{Computational Overhead}
Soft-NMS adds $<$0.3ms per frame over Hard NMS (from 32 FPS to
29 FPS). This 9.4\% throughput reduction is negligible compared to
the 57\% improvement in counting accuracy. The density head adds
only 0.2ms per forward pass, representing $<$1\% of total inference
time.

\subsection{Limitations}
\begin{itemize}
  \item The current pipeline assumes axis-aligned bounding boxes.
    Rotated or irregular objects (biological cells, angled products)
    would benefit from oriented bounding box regression.
  \item $\sigma$ is a global parameter; per-region adaptive $\sigma$
    (inspired by Adaptive NMS) could further improve performance
    in scenes with heterogeneous density.
  \item The robustness analysis covers additive perturbations but
    does not test adversarial or geometric distortions.
\end{itemize}

% ============================================================
\section{Conclusion \& Future Work}
\label{sec:conclusion}

We presented a comprehensive study of instance-level detection for
densely overlapping objects, progressing from heuristic baselines through
classical ML to deep learning with Soft-NMS. Our key findings are:

\begin{enumerate}
  \item Soft-NMS with $\sigma = 0.5$ yields consistent improvements
    across all density levels, achieving 81.7\% mAP@0.5 and an MAE
    of 1.8 objects---a 57\% error reduction over Hard NMS.
  \item The density estimation head provides effective regularization
    during training and fast approximate counting during inference.
  \item Performance degrades gracefully under noise, blur, and
    brightness perturbations, with worst-case mAP reduction of 4.1\%.
  \item The Random Forest ML baseline validates the hybrid classical
    pipeline concept while highlighting the representation gap between
    hand-crafted and learned features.
\end{enumerate}

\textbf{Future work} includes: (a) extending to pixel-level instance
masks via YOLACT-style prototype generation; (b) exploring per-region
adaptive $\sigma$ for heterogeneous density scenes; (c) INT8
quantization for edge TPU deployment; and (d) scaling synthetic data
generation with more complex occlusion patterns.

% ============================================================
\bibliographystyle{IEEEtran}
\bibliography{references}

\end{document}
